% GENERATED FILE. Edit canonical lessons, then run npm run generate:books.
% canonical-source-hash: fnv1a64:f36dfe86e77e35a9
% canonical-lessons: ZH-C270-tongyi, ZH-C270-fandui, ZH-C270-zancheng, ZH-C270-zhichi, ZH-C270-xiangxin

\chapter{Do You Agree?}
\label{ch:zh-do-you-agree}
\hlchaptermodality{270}

\begin{chapteropening}
\textbf{By the end of this chapter:} \emph{I can say I agree or that I am against something, that I am in favour, that I am behind someone, and that I believe them.}

Complete the last lesson of chapter 270: I can say I agree or that I am against something, that I am in favour, that I am behind someone, and that I believe them.
\end{chapteropening}

\section[tóngyì]{tóngyì — to agree}

\label{lesson:ZH-C270-tongyi}



\begin{quote}
Before the new one: say the Mandarin for originally, then the Mandarin for a toy.
\end{quote}

\subsection*{You'll want to know: tóngyì}
\emph{tóngyì} — \textquotedblleft{}to agree\textquotedblright{}.

\begin{quote}
\emph{nǐ tóngyì ma} — Do you agree?
\end{quote}

Its characters come later; for now it is a word for the ear.

\begin{cousinweb}
\emph{tóngyì} is built from \zh{同} 'the same' and \emph{yì} 'a mind', both inherited from Middle Chinese, \emph{duwng} and \emph{'iH} (Guangyun, in Baxter's transcription), which glosses the second 'will': to be of one mind.
\end{cousinweb}

\begin{grammarlens}[title={What you've built}]
One more word for saying no and asking.
\end{grammarlens}

\subsection*{Your turn}
\begin{itemize}
\raggedright
  \item \emph{Say it:} \emph{tóngyì}
  \item \emph{Say it:} \emph{tóngyì}, once more
  \item \emph{Say it:} \emph{nǐ tóngyì ma} — Do you agree?
  \item \emph{Recall:} say the Mandarin for originally, then the Mandarin for a toy, then say \emph{tóngyì} again
\end{itemize}

\begin{tcolorbox}[breakable,colback=teal!4,colframe=teal!35!black,title={Before you move on}]
What does tóngyì mean? (\textquotedblleft{}To agree\textquotedblright{}.) Say it once more.
\end{tcolorbox}
\section[fǎnduì]{fǎnduì — to be against; to oppose}

\label{lesson:ZH-C270-fandui}



\begin{quote}
Before the new one: say the Mandarin for a toy, then the Mandarin for to agree.
\end{quote}

\subsection*{You'll want to know: fǎnduì}
\emph{fǎnduì} — \textquotedblleft{}to be against; to oppose\textquotedblright{}.

\begin{quote}
\emph{wǒ bù fǎnduì} — I'm not against it.
\end{quote}

Its characters come later; for now it is a word for the ear.

\begin{cousinweb}
\emph{fǎnduì} is built from \emph{fǎn} and \zh{对}, both inherited from Middle Chinese, \emph{pjonX} and \emph{twojH} (Guangyun, in Baxter's transcription), which glosses the first 'to turn over' and 'contrary', and the second 'to answer'.
\end{cousinweb}

\begin{grammarlens}[title={What you've built}]
One more word for saying no and asking.
\end{grammarlens}

\subsection*{Your turn}
\begin{itemize}
\raggedright
  \item \emph{Say it:} \emph{fǎnduì}
  \item \emph{Say it:} \emph{fǎnduì}, once more
  \item \emph{Say it:} \emph{wǒ bù fǎnduì} — I'm not against it.
  \item \emph{Recall:} say the Mandarin for a toy, then the Mandarin for to agree, then say \emph{fǎnduì} again
\end{itemize}

\begin{tcolorbox}[breakable,colback=teal!4,colframe=teal!35!black,title={Before you move on}]
What does fǎnduì mean? (\textquotedblleft{}To be against; to oppose\textquotedblright{}.) Say it once more.
\end{tcolorbox}
\section[zànchéng]{zànchéng — to approve of; to be in favour}

\label{lesson:ZH-C270-zancheng}



\begin{quote}
Before the new one: say the Mandarin for to agree, then the Mandarin for to be against.
\end{quote}

\subsection*{You'll want to know: zànchéng}
\emph{zànchéng} — \textquotedblleft{}to approve of; to be in favour\textquotedblright{}.

\begin{quote}
\emph{dàjiā dōu zànchéng} — Everyone's in favour.
\end{quote}

Its characters come later; for now it is a word for the ear.

\begin{cousinweb}
\emph{zànchéng} is built from \emph{zàn} 'to assist' and \emph{chéng} 'to complete', both inherited from Middle Chinese, \emph{tsanH} and \emph{dzyeng} (Guangyun, in Baxter's transcription), which glosses the first 'to assist' and the second 'to finish'.
\end{cousinweb}

\begin{grammarlens}[title={What you've built}]
One more word for saying no and asking.
\end{grammarlens}

\subsection*{Your turn}
\begin{itemize}
\raggedright
  \item \emph{Say it:} \emph{zànchéng}
  \item \emph{Say it:} \emph{zànchéng}, once more
  \item \emph{Say it:} \emph{dàjiā dōu zànchéng} — Everyone's in favour.
  \item \emph{Recall:} say the Mandarin for to agree, then the Mandarin for to be against, then say \emph{zànchéng} again
\end{itemize}

\begin{tcolorbox}[breakable,colback=teal!4,colframe=teal!35!black,title={Before you move on}]
What does zànchéng mean? (\textquotedblleft{}To approve of; to be in favour\textquotedblright{}.) Say it once more.
\end{tcolorbox}
\section[zhīchí]{zhīchí — to support; to back}

\label{lesson:ZH-C270-zhichi}



\begin{quote}
Before the new one: say the Mandarin for to be against, then the Mandarin for to be in favour.
\end{quote}

\subsection*{You'll want to know: zhīchí}
\emph{zhīchí} — \textquotedblleft{}to support; to back\textquotedblright{}.

\begin{quote}
\emph{wǒ zhīchí nǐ} — I'm behind you.
\end{quote}

Its characters come later; for now it is a word for the ear.

\begin{cousinweb}
\emph{zhīchí} is built from \emph{zhī} and \emph{chí} 'to hold', both inherited from Middle Chinese, \emph{tsye} and \emph{dri} (Guangyun, in Baxter's transcription), which glosses the first with this very word and the second 'to hold'.
\end{cousinweb}

\begin{grammarlens}[title={What you've built}]
One more word for saying no and asking.
\end{grammarlens}

\subsection*{Your turn}
\begin{itemize}
\raggedright
  \item \emph{Say it:} \emph{zhīchí}
  \item \emph{Say it:} \emph{zhīchí}, once more
  \item \emph{Say it:} \emph{wǒ zhīchí nǐ} — I'm behind you.
  \item \emph{Recall:} say the Mandarin for to be against, then the Mandarin for to be in favour, then say \emph{zhīchí} again
\end{itemize}

\begin{tcolorbox}[breakable,colback=teal!4,colframe=teal!35!black,title={Before you move on}]
What does zhīchí mean? (\textquotedblleft{}To support; to back\textquotedblright{}.) Say it once more.
\end{tcolorbox}
\section[xiāngxìn]{xiāngxìn — to believe; to trust}

\label{lesson:ZH-C270-xiangxin}



\begin{quote}
Before the new one: say the Mandarin for to be in favour, then the Mandarin for to back someone up.
\end{quote}

\subsection*{You'll want to know: xiāngxìn}
\emph{xiāngxìn} — \textquotedblleft{}to believe; to trust\textquotedblright{}.

\begin{quote}
\emph{wǒ xiāngxìn nǐ} — I believe you.
\end{quote}

Its characters come later; for now it is a word for the ear.

\begin{cousinweb}
\emph{xiāngxìn} is built from \emph{xiāng} 'each other' and \emph{xìn} 'to trust', both inherited from Middle Chinese, \emph{sjang} and \emph{sinH} (Guangyun, in Baxter's transcription), which glosses the first 'together' and the second 'faithful'.
\end{cousinweb}

\begin{grammarlens}[title={What you've built}]
One more word for saying no and asking.
\end{grammarlens}

\subsection*{Your turn}
\begin{itemize}
\raggedright
  \item \emph{Say it:} \emph{xiāngxìn}
  \item \emph{Say it:} \emph{xiāngxìn}, once more
  \item \emph{Say it:} \emph{wǒ xiāngxìn nǐ} — I believe you.
  \item \emph{Recall:} say the Mandarin for to be in favour, then the Mandarin for to back someone up, then say \emph{xiāngxìn} again
\end{itemize}

\begin{tcolorbox}[breakable,colback=teal!4,colframe=teal!35!black,title={Before you move on}]
What does xiāngxìn mean? (\textquotedblleft{}To believe; to trust\textquotedblright{}.) Say it once more.
\end{tcolorbox}
